\section*{第 \ref{ch:g10:algebra} 章 --- 代数：方程与不等式}

\begin{solution}{exo:g10:algebra:1}
$3(2x-5) - 2(x-7) = 6x - 15 - 2x + 14 = 4x - 1$.

$(x+4)(2x-1) = 2x^2 - x + 8x - 4 = 2x^2 + 7x - 4$.

$(3x-2)^2 = 9x^2 - 12x + 4$.

$(5-x)(5+x) = 25 - x^2$.
\end{solution}

\begin{solution}{exo:g10:algebra:2}
$x^2 - 49 = (x+7)(x-7)$（平方差）。

$x^2 + 10x + 25 = (x+5)^2$.

$7x^2 - 14x = 7x(x - 2)$（提取公因式）。

$16x^2 - 8x + 1 = (4x - 1)^2$.
\end{solution}

\begin{solution}{exo:g10:algebra:3}
$4x + 9 = x - 3$：$3x = -12$，故 $x = -4$。

$\frac{2x-1}{3} = 5$：$2x - 1 = 15$，故 $x = 8$。

$2(x-3) = 2x + 1$ 展开得 $2x - 6 = 2x + 1$，即 $-6 = 1$，
对任意 $x$ 都不成立。无解——两边分别表示的两条直线
平行，永不相交。
\end{solution}

\begin{solution}{exo:g10:algebra:4}
$(x+5)(2x-8) = 0$：$x = -5$ 或 $x = 4$。

$x^2 - 16 = (x-4)(x+4) = 0$：$x = 4$ 或 $x = -4$。

$x^2 + 6x + 9 = (x+3)^2 = 0$：唯一解为 $x = -3$。
\end{solution}

\begin{solution}{exo:g10:algebra:5}
$3x - 5 \leq 7$：$3x \leq 12$，故 $x \leq 4$，解集为
$\intoc{-\infty}{4}$。

$-4x + 1 > 9$：$-4x > 8$，两边除以 $-4$，不等号方向改变：
$x < -2$，解集为 $\intoo{-\infty}{-2}$。

$2(x-1) \geq 5x + 4$：$2x - 2 \geq 5x + 4$，故 $-3x \geq 6$，于是
$x \leq -2$，解集为 $\intoc{-\infty}{-2}$。
\end{solution}

\begin{solution}{exo:g10:algebra:6}
公因式为 $(x-1)$：
\[
E = (x-1)\bigl[(x+3) - (2x - 5)\bigr] = (x-1)(x + 3 - 2x + 5)
= (x-1)(8 - x).
\]
于是，$E = 0$ 当 $x = 1$ 或 $x = 8$ 时成立（零乘积法则）。
\end{solution}

\begin{solution}{exo:g10:algebra:7}
$x^2 = 5x$：$x^2 - 5x = x(x - 5) = 0$，故 $x = 0$ 或 $x = 5$（切勿
两边同除以 $x$，否则会丢失解 $0$）。

$(2x+1)^2 = (x-4)^2$：移项，并利用平方差
公式分解：
\[
(2x+1)^2 - (x-4)^2
= \bigl[(2x+1) + (x-4)\bigr]\bigl[(2x+1) - (x-4)\bigr]
= (3x - 3)(x + 5).
\]
由零乘积法则得 $x = 1$ 或 $x = -5$。
\end{solution}

\begin{solution}{exo:g10:algebra:8}
$(x+2)(4-x) > 0$：因式 $x + 2$ 在 $-2$ 处为零（其后为正），
$4 - x$ 在 $4$ 处为零（其前为正）。乘积为正
当且仅当两个因式同时为正，即在 $\intoo{-2}{4}$ 上。

$\frac{2x-6}{x+1} \geq 0$：分子在 $3$ 处为零，分母在
$-1$ 处为零（不允许取值）。符号表：商为正当且仅当两部分
同号，即 $x < -1$ 或 $x > 3$；商在 $x = 3$ 处为零。
解集为 $\intoo{-\infty}{-1} \cup \intco{3}{+\infty}$。
\end{solution}

\begin{solution}{exo:g10:algebra:9}
经过 $x$ 个月后，第一种服务的费用为 $9x + 15$，第二种为 $12x$。
第一种更便宜的条件是
\[
9x + 15 < 12x
\ \Longleftrightarrow\
15 < 3x
\ \Longleftrightarrow\
x > 5 .
\]
从第 $6$ 个月起，第一种服务总体上是更划算的选择。
\end{solution}

\begin{solution}{exo:g10:algebra:10}
不允许的值：$x = 1$。当 $x \neq 1$ 时，两边同乘 $x - 1$：
\[
x + 3 = 2(x - 1) = 2x - 2
\ \Longleftrightarrow\
x = 5 .
\]
由于 $5 \neq 1$，故解为 $x = 5$。验算：
$\frac{5+3}{5-1} = \frac84 = 2$.
\end{solution}

\begin{solution}{exo:g10:algebra:11}
$(n+1)^2 - n^2 = n^2 + 2n + 1 - n^2 = 2n + 1$，它对
每个整数 $n$ 都是奇数。

两个相邻奇数可写成 $2n - 1$ 和 $2n + 1$。于是
\[
(2n+1)^2 - (2n-1)^2
= \bigl[(2n+1)+(2n-1)\bigr]\bigl[(2n+1)-(2n-1)\bigr]
= 4n \times 2 = 8n,
\]
为 $8$ 的倍数。
\end{solution}

\begin{solution}{pb:g10:algebra:1}
\textbf{1.} $x^2 - 10x + 25 = (x - 5)^2$；
$9x^2 - 49 = (3x - 7)(3x + 7)$；
$x^2 + 6x + 9 - 4 = (x + 3)^2 - 2^2 = (x + 1)(x + 5)$。

\textbf{2.} $x = -1$ 或 $x = -5$；\quad $x = \frac73$ 或
$x = -\frac73$。

\textbf{3.} 由 $(5 - t)(5 + t) = 25 - t^2 = 21$ 得 $t^2 = 4$，
$t = 2$：所求两数为 $3$ 和 $7$（和为 $10$，积为 $21$，
验算）。

\textbf{4.} $(7 - t)(7 + t) = 49 - t^2 = 33$：$t = 4$，两数
为 $3$ 和 $11$。

\textbf{5.} $\left(\frac s2 - t\right)\left(\frac s2 + t\right)
= \left(\frac s2\right)^2 - t^2 = p$ 可推出
$t^2 = \left(\frac s2\right)^2 - p$。方程有解当且仅当
此式 $\geq 0$，即
$p \leq \left(\frac s2\right)^2$ ——和为定值的两数，其积
永远不会超过半和的平方
（狄多女王的围栏问题，初中分册中已见过）。

\textbf{6.} $(x - a)(x - b) = x^2 - (a + b)x + ab$：方程
$x^2 - sx + p = 0$ 恰好表明“$x$ 是两数之一，
这两数的和为 $s$、积为 $p$”。

\textbf{7.} $(x-5)^2 - 4 = x^2 - 10x + 25 - 4 = x^2 - 10x +
21$。应用平方差公式：
$(x - 5 - 2)(x - 5 + 2) = (x - 7)(x - 3) = 0$：解为 $3$
和 $7$ ——正是第 3 题中书吏所求的两数。

\textbf{8.} $x^2 + 6x - 7 = (x + 3)^2 - 16 = 0$ 给出
$(x + 3)^2 = 16$，$x + 3 = \pm 4$：$x = 1$ 或 $x = -7$。

\textbf{9.} $x^2 + 4x + 5 = (x + 2)^2 + 1 \geq 1 > 0$ 对每个
$x$ 都成立，故无实数解。写成 $(x + h)^2 + k$ 的形式时，
当且仅当 $k \leq 0$ 时有解（此时 $(x+h)^2 = -k$ 的两根
为 $-h \pm \sqrt{-k}$）；当 $k > 0$ 时，
该式恒为正。

\textbf{10.} 边长为 $x$ 的正方形（面积 $x^2$）与两个
$3 \times x$ 的矩形（共 $6x$）构成一个 L 形，其面积为
$x^2 + 6x = 7$；缺失的 $3 \times 3$ 角块（面积 $9$）
恰好补成边长为 $x + 3$ 的大正方形。两边同加 $9$
得 $(x + 3)^2 = 16$，故 $x + 3 = 4$（长度为正）
于是 $x = 1$。代数运算与图形所做的其实是同一件事。

\textbf{11.} $(x-3)(x+2)$：两根之间为负。解
$< 0$ 的解：$x \in \intoo{-2}{3}$。

\textbf{12.} $x^2 - 10x + 21 \leq 0$，即
$(x - 3)(x - 7) \leq 0$：两根之间（含端点）：
$x \in \intcc{3}{7}$。

\textbf{13.} 零点为 $x = 1$，$x = -2$ 不允许取值；
分式在两根之外为正：
$x \in \intoo{-\infty}{-2} \cup \intco{1}{+\infty}$。

\textbf{14.} $25 - (x - 5)^2 = 25 - x^2 + 10x - 25 =
x(10 - x)$：恒等式已验证。由于 $(x-5)^2 \geq 0$ 恒成立，
$A(x) \leq 25$，等号成立当且仅当 $x = 5$：
$5 \times 5$ 的正方形围栏，最大值 $25$ m$^2$ ——六年级
见过的表格，如今只需两行即可证明。

\textbf{15.} 设 $x > 0$，将命题两边乘以 $x$：
$x^2 + 1 \geq 2x$，即 $x^2 - 2x + 1 = (x - 1)^2 \geq 0$ ——
成立，仅当 $x = 1$ 时取等号。再除以 $x > 0$
结论不变（见\cref{prop:g10:algebra:ineqrules}）。

\textbf{16.} 对 $\frac1x = \frac{x}{1 - x}$ 交叉相乘
（各长度均为正）：$x^2 = 1 - x$，即
$x^2 + x - 1 = 0$。

\textbf{17.} $x^2 + x - 1 = \left(x + \frac12\right)^2 -
\frac54 = 0$ 得 $x + \frac12 = \pm\frac{\sqrt5}{2}$。
取正根：$x = \frac{\sqrt5 - 1}{2} \approx 0.618$。

\textbf{18.} $\varphi = \frac{2}{\sqrt5 - 1} =
\frac{2(\sqrt5 + 1)}{(\sqrt5 - 1)(\sqrt5 + 1)}
= \frac{2(\sqrt5 + 1)}{5 - 1} = \frac{\sqrt5 + 1}{2}
\approx 1.618$。

\textbf{19.} $\varphi^2 = \frac{(\sqrt5 + 1)^2}{4} =
\frac{6 + 2\sqrt5}{4} = \frac{3 + \sqrt5}{2}$，且
$\varphi + 1 = \frac{\sqrt5 + 3}{2}$：两者相等。又
$\varphi - 1 = \frac{\sqrt5 - 1}{2} = x = \frac1\varphi$（这
正是 $\varphi$ 的定义方式）。黄金比是这样的数：
其平方等于它加 1，其倒数等于它减 1。

\textbf{20.} $\sqrt5$ 是无理数（初中分册中关于无理数的周末问题：$5$
不是完全平方数），加 $1$ 再减半后
仍为无理数（有理数的运算，见\cref{pb:g10:numbers:1}）。
比值：$\frac53 \approx 1.667$，$\frac85 = 1.600$，
$\frac{13}{8} = 1.625$，$\frac{21}{13} \approx 1.615$：
交替位于其上下，逐步逼近
$\varphi = 1.618\ldots$ ——初中分册“节奏计数”周末问题的数，
暗中追逐着黄金分割；
数列相关章节将当场捉住它们。
\end{solution}
